% GENERATED FILE. Edit canonical lessons, then run npm run generate:books.
% canonical-source-hash: fnv1a64:9d1c5014b181d11c
% canonical-lessons: PT-C519-ensaio, PT-C519-ensaiar, PT-C519-palco, PT-C519-plateia, PT-C519-aplauso

\chapter{Antes do espetáculo}
\label{ch:pt-antes-do-espetaculo}
\hlchaptermodality{519}

\begin{chapteropening}
\textbf{By the end of this chapter:} \emph{I can talk about rehearsals, the stage and the audience.}

Complete the last lesson of chapter 519: I can talk about rehearsals, the stage and the audience.
\end{chapteropening}

\section[o ensaio]{o ensaio — rehearsal; essay}

\label{lesson:PT-C519-ensaio}



\begin{quote}
Before the new one: say the Portuguese for to hum, then the Portuguese for to whistle.
\end{quote}

\subsection*{You'll want to know: o ensaio}
\textbf{o ensaio} — \textquotedblleft{}rehearsal; essay\textquotedblright{}.

\begin{quote}
O ensaio da peça é às oito.
\end{quote}

\begin{quote}
\emph{The rehearsal for the play is at eight.}
\end{quote}

Also a written essay: \emph{um ensaio sobre arte}.

\begin{cousinweb}
\emph{ensaio} is inherited, through Old Portuguese \emph{eisaio}, from Late Latin \emph{exagium} \textquotedblleft{}a weighing, a trial\textquotedblright{} (EtymDB), related to \emph{exigere} \textquotedblleft{}to weigh out\textquotedblright{}.
\end{cousinweb}

\begin{grammarlens}[title={What you've built}]
One more word for telling, explaining and describing at B1.
\end{grammarlens}

\subsection*{Your turn}
\begin{itemize}
\raggedright
  \item \emph{Say it:} \emph{o ensaio}
  \item \emph{Say it:} \emph{o ensaio}, once more
  \item \emph{Say it:} a sentence of your own with \emph{o ensaio}
  \item \emph{Recall:} say the Portuguese for to hum, then the Portuguese for to whistle, then say \emph{o ensaio} again
\end{itemize}

\begin{tcolorbox}[breakable,colback=teal!4,colframe=teal!35!black,title={Before you move on}]
What does o ensaio mean? (\textquotedblleft{}Rehearsal; essay\textquotedblright{}.) Say it once more.
\end{tcolorbox}
\section[ensaiar]{ensaiar — to rehearse}

\label{lesson:PT-C519-ensaiar}



\begin{quote}
Before the new one: say the Portuguese for to whistle, then the Portuguese for rehearsal.
\end{quote}

\subsection*{You'll want to know: ensaiar}
\textbf{ensaiar} — \textquotedblleft{}to rehearse\textquotedblright{}.

\begin{quote}
A banda ensaia na garagem.
\end{quote}

\begin{quote}
\emph{The band rehearses in the garage.}
\end{quote}

Built on \emph{o ensaio}.

\begin{cousinweb}
\emph{ensaiar} is built on \emph{ensaio} \textquotedblleft{}a trial, a rehearsal\textquotedblright{}, inherited from Late Latin \emph{exagium} \textquotedblleft{}a weighing\textquotedblright{} (EtymDB).
\end{cousinweb}

\begin{grammarlens}[title={What you've built}]
One more word for telling, explaining and describing at B1.
\end{grammarlens}

\subsection*{Your turn}
\begin{itemize}
\raggedright
  \item \emph{Say it:} \emph{ensaiar}
  \item \emph{Say it:} \emph{ensaiar}, once more
  \item \emph{Say it:} a sentence of your own with \emph{ensaiar}
  \item \emph{Recall:} say the Portuguese for to whistle, then the Portuguese for rehearsal, then say \emph{ensaiar} again
\end{itemize}

\begin{tcolorbox}[breakable,colback=teal!4,colframe=teal!35!black,title={Before you move on}]
What does ensaiar mean? (\textquotedblleft{}To rehearse\textquotedblright{}.) Say it once more.
\end{tcolorbox}
\section[o palco]{o palco — stage}

\label{lesson:PT-C519-palco}



\begin{quote}
Before the new one: say the Portuguese for rehearsal, then the Portuguese for to rehearse.
\end{quote}

\subsection*{You'll want to know: o palco}
\textbf{o palco} — \textquotedblleft{}stage\textquotedblright{}.

\begin{quote}
Os atores entraram no palco.
\end{quote}

\begin{quote}
\emph{The actors came on stage.}
\end{quote}

\emph{Subir ao palco} — to go on stage.

\begin{cousinweb}
\emph{palco} was borrowed from Italian \emph{palco} \textquotedblleft{}a platform, a stage\textquotedblright{} (EtymDB), a word of Germanic origin.
\end{cousinweb}

\begin{grammarlens}[title={What you've built}]
One more word for telling, explaining and describing at B1.
\end{grammarlens}

\subsection*{Your turn}
\begin{itemize}
\raggedright
  \item \emph{Say it:} \emph{o palco}
  \item \emph{Say it:} \emph{o palco}, once more
  \item \emph{Say it:} a sentence of your own with \emph{o palco}
  \item \emph{Recall:} say the Portuguese for rehearsal, then the Portuguese for to rehearse, then say \emph{o palco} again
\end{itemize}

\begin{tcolorbox}[breakable,colback=teal!4,colframe=teal!35!black,title={Before you move on}]
What does o palco mean? (\textquotedblleft{}Stage\textquotedblright{}.) Say it once more.
\end{tcolorbox}
\section[a plateia]{a plateia — audience; the stalls}

\label{lesson:PT-C519-plateia}



\begin{quote}
Before the new one: say the Portuguese for to rehearse, then the Portuguese for stage.
\end{quote}

\subsection*{You'll want to know: a plateia}
\textbf{a plateia} — \textquotedblleft{}audience; the stalls\textquotedblright{}.

\begin{quote}
A plateia estava cheia.
\end{quote}

\begin{quote}
\emph{The auditorium was full.}
\end{quote}

Also the seats on the ground floor of a theatre.

\begin{cousinweb}
\emph{plateia} comes from Latin \emph{platēa} \textquotedblleft{}a broad street, a courtyard\textquotedblright{} (EtymDB), from Greek \emph{plateîa} \textquotedblleft{}broad\textquotedblright{}.
\end{cousinweb}

\begin{grammarlens}[title={What you've built}]
One more word for telling, explaining and describing at B1.
\end{grammarlens}

\subsection*{Your turn}
\begin{itemize}
\raggedright
  \item \emph{Say it:} \emph{a plateia}
  \item \emph{Say it:} \emph{a plateia}, once more
  \item \emph{Say it:} a sentence of your own with \emph{a plateia}
  \item \emph{Recall:} say the Portuguese for to rehearse, then the Portuguese for stage, then say \emph{a plateia} again
\end{itemize}

\begin{tcolorbox}[breakable,colback=teal!4,colframe=teal!35!black,title={Before you move on}]
What does a plateia mean? (\textquotedblleft{}Audience; the stalls\textquotedblright{}.) Say it once more.
\end{tcolorbox}
\section[o aplauso]{o aplauso — applause}

\label{lesson:PT-C519-aplauso}



\begin{quote}
Before the new one: say the Portuguese for stage, then the Portuguese for audience.
\end{quote}

\subsection*{You'll want to know: o aplauso}
\textbf{o aplauso} — \textquotedblleft{}applause\textquotedblright{}.

\begin{quote}
No fim, os músicos receberam um grande aplauso.
\end{quote}

\begin{quote}
\emph{At the end, the musicians got a big round of applause.}
\end{quote}

Often plural: \emph{os aplausos}.

\begin{cousinweb}
\emph{aplauso} comes from Latin \emph{applausus} \textquotedblleft{}applause\textquotedblright{} (EtymDB), from \emph{applaudere}, built on \emph{plaudere} \textquotedblleft{}to clap\textquotedblright{}.
\end{cousinweb}

\begin{grammarlens}[title={What you've built}]
One more word for telling, explaining and describing at B1.
\end{grammarlens}

\subsection*{Your turn}
\begin{itemize}
\raggedright
  \item \emph{Say it:} \emph{o aplauso}
  \item \emph{Say it:} \emph{o aplauso}, once more
  \item \emph{Say it:} a sentence of your own with \emph{o aplauso}
  \item \emph{Recall:} say the Portuguese for stage, then the Portuguese for audience, then say \emph{o aplauso} again
\end{itemize}

\begin{tcolorbox}[breakable,colback=teal!4,colframe=teal!35!black,title={Before you move on}]
What does o aplauso mean? (\textquotedblleft{}Applause\textquotedblright{}.) Say it once more.
\end{tcolorbox}
